% Replaces the lumen volume to myocardial mass section (2026-11-02, on request); V/M stays a
% secondary candidate in 04_model_size.tex (sec:method-selection).
% Implementation: tortuosity/taper.py (scored by tortuosity/score.py). Not yet in
% CGPS_data/thesis_final/. Origin of the idea: reports/Fourth centerline feature for plaque.md;
% physics and branch flow loss: reports/Taper and curvature interaction for WSS.md.
% Background: the taper studies (johnston2009taper, schultz2023ess) are described in the literature
% review (sec:lit-wss, 2026-11-02); the physics rationale is the first paragraph here.
\providecommand{\clarify}[1]{\textcolor{red}{\textbf{Clarification needed:} #1}}
\subsection{Lumen Taper}
\label{sec:method-taper}

Taper was included for its place in the relation between lumen size and wall shear stress, not for direct evidence that it predicts plaque.
For steady flow through a straight tube, wall shear stress equals $4\mu Q/(\pi r^3)$, where $\mu$ is the blood viscosity, $Q$ the flow and $r$ the lumen radius \cite{malek1999shear}.
At a given flow, a segment that holds its calibre or widens therefore has lower wall shear stress than one that narrows.
Along a real artery, however, flow leaves through every side branch, so $Q$ falls together with $r$.
If flow scaled with $r^3$, as Murray's law predicts, wall shear stress would stay constant along the vessel and taper would carry no information on it; in coronary trees the pooled exponent is 2.39 \cite{taylor2024murray}, so taper is expected to carry only part of the signal that the tube formula suggests.

The evidence that taper itself matters is thin and comes from idealised simulations and one study of normal arteries, and no study has related taper to later plaque (Section~\ref{sec:lit-wss}).
Taper was nevertheless chosen as the fourth feature, as lumen size is the geometric quantity that the wall shear stress formula weights most, and a slope of the logarithm of the radius describes how the size changes along the vessel without depending on the overall calibre, which is set by sex and body size.
A lower taper is expected to mean lower wall shear stress and more plaque, whereas diffuse wall thickening below the detection limit of CCTA could steepen taper before plaque is graded and give the opposite sign.
\clarify{whether the test of taper is two-sided or carries a prespecified sign (plan).}

Normal coronary arteries lose about 0.25\,mm of diameter per 10\,mm of length \cite{medranogracia2016atlas}, which over a 30\,mm stretch is a change in radius of about one voxel.
No single cross-section resolves this, so taper was estimated by a robust fit over many cross-sections.
The centerline of each artery (Section~\ref{sec:method-split}) was resampled every 0.5\,mm of arc length $s$, measured from the start of the artery.
At each sample, the tangent was taken from the points 2.5\,mm before and after it, cut short at the ends of the artery, so that taper and curvature (Section~\ref{sec:method-curvature}) are smoothed over the same 5\,mm.
The lumen was cut with the plane orthogonal to the tangent on a grid of $12 \times 12$\,mm with 0.1\,mm spacing, reading the lumen mask with trilinear interpolation and a threshold of 0.5 so that the area resolves below one voxel.
The cross-sectional area $A$ was that of the connected region containing the centerline point, and the radius was the area-equivalent radius
\begin{equation}
  r(s) = \sqrt{A(s)/\pi}.
  \label{eq:method-taper-radius}
\end{equation}
The area-equivalent radius was used rather than the distance-transform radius, which follows the minor axis of a non-circular section; the distance-transform radius was kept for a sensitivity analysis.
\clarify{the code reads a voxel lumen mask; state how the mask is formed from the tracer's lumen contours in the CGPS scans.}

Samples were excluded in the first 3\,mm of the artery and within two median radii of a junction with a side branch, about one parent diameter, where the radius jumps at the branch origin.
Samples were also excluded where the centerline point lay outside the lumen, and where the connected region reached the edge of the grid or contained a voxel of another labelled vessel.
The taper was then
\begin{equation}
  T = -10\,b,
  \label{eq:method-taper}
\end{equation}
where $b$ is the Theil--Sen slope of $\ln r$ against $s$ over the remaining samples, in mm$^{-1}$, so that $T$ is the fractional loss of radius per 10\,mm.
A positive $T$ means normal tapering, and a $T$ near zero or negative means an artery that holds its calibre or widens.
The Theil--Sen estimator, the median of the slopes between all pairs of samples, was used so that remaining steps at small side branches and local errors in the contour do not drive the fit.
The taper of an artery was recorded as missing when the remaining samples spanned less than 20\,mm or made up less than 60\% of all samples.
\clarify{tortuosity/taper.py computes one $T$ per artery (the whole LAD, LCx or RCA), whereas Section~\ref{sec:method-selection} uses taper per segment; either the code is run per segment, or the text and the model use one value per artery.}
